\section{Luttinger--Ward 泛函与骨架展开}

把相互作用贡献写成 $G$ 的泛函 $\Phi[G]$（Luttinger--Ward），使得
\begin{equation}
  \Gamma[G]
  =\Gamma_0[G]+\Phi[G],
  \qquad
  \Sigma[G]
  =\frac{\delta\Phi}{\delta G}.
\label{eq:LW}
\end{equation}
$\Phi$ 由\textbf{骨架图}（用全 $G$ 线、无自能插入）构成。由此自动得到 Dyson 方程
\begin{equation}
  G^{-1}=G_0^{-1}-\Sigma[G].
\end{equation}
Baym--Kadanoff 守恒近似即要求 $\Sigma$ 由某个 $\Phi$ 导出，从而保证粒子数、动量、能量等守恒律。

常见近似：
\begin{itemize}
  \item HF：$\Phi$ 取一阶直接+交换；
  \item GW：$\Phi$ 取环图级数（屏蔽相互作用 $W$）；
  \item 涨落交换（FLEX）等更高阶梯子/气泡组合。
\end{itemize}

二粒子不可约顶点
\begin{equation}
  \Gamma^{(4)}
  =\frac{\delta^2\Phi}{\delta G\,\delta G}
\end{equation}
进入 Bethe--Salpeter 方程，决定响应函数与集体模式。
